%======================================================================%
\section{The Induced Yang--Mills Coefficient}
\label{sec:app-ym:main}
%======================================================================%

This appendix derives the one-loop coefficient of the composite
$\SOTen\times\UOne$ Yang--Mills term obtained by integrating out the
massive clock fields, and traces how that coefficient feeds the
asymptotic-freedom slope $b_{0}$ of the functional renormalization group.
The computation is presented at \emph{derivation grade}: the heat-kernel
machinery, the group-theory weights, and the resulting logarithm are all
theorems given the stated background; the embedding of the clock-field
$\sigma$-model into a four-dimensional Lorentzian arena is the soldering
\emph{postulate} P3, invoked only to provide the $X=M^{4}$ over which the
proper-time integral is taken.

Throughout we use the metric signature $(-,+,+,+)$, the decay constant
$f$ with $[f]=\text{mass}$, the breaking/mass scale $\mu$, and the UV
cutoff $\Lambda$. The hatted coupling $\ghat^{2}$ is dimensionless. The
relevant gauge algebra is $\mathfrak{so}_{10}\oplus\mathfrak u_{1}$, and
\emph{all} group-theory weights are taken in $\SOTen$ with the locked
Casimirs
\begin{equation}\label{eq:app-ym:casimirs}
  C_{A}(\SOTen)=8,\qquad
  T(\mathbf{16})=2,\qquad
  C_{2}(\mathbf{16})=\tfrac{45}{8}.
\end{equation}
The $E_{6}$ value $C_{A}(E_{6})=12$ is \emph{never} mixed into these
sums: it is the spinor index $T(\mathbf{16})=2$, not a count of sixteen
index-$1$ fundamentals, that enters.

%----------------------------------------------------------------------%
\subsection{The Maurer--Cartan curvature and the one-loop determinant}
\label{sec:app-ym:setup}
%----------------------------------------------------------------------%

Let $g(x)\in E_{6}$ (compact real form, the dynamical form fixed in
\autoref{sec:foundation:realform}) be a coset representative of the
vacuum map into $M=E_{6}/(\SOTen\times\UOne)=\mathrm{EIII}$. The
Maurer--Cartan form splits along the reductive decomposition
$\mathfrak e_{6}=\mathfrak h\oplus\mathfrak m$,
\begin{equation}\label{eq:app-ym:MC}
  g^{-1}\partial_{\mu}g
  = \mathcal A_{\mu}(x) + e_{\mu}(x),
  \qquad
  \mathcal A_{\mu}\in\mathfrak{so}_{10}\oplus\mathfrak u_{1},
  \qquad
  e_{\mu}\in\mathfrak m
  =\mathbf{16}_{-3}\oplus\overline{\mathbf{16}}_{+3},
\end{equation}
with $\mathfrak h$-part $\mathcal A_{\mu}$ the composite connection of
\autoref{sec:gauge:composite} and $\mathfrak m$-part $e_{\mu}$ the
clock-field vielbein. The $\mathfrak h$-curvature
\begin{equation}\label{eq:app-ym:curv}
  \mathcal F_{\mu\nu}
  =\partial_{\mu}\mathcal A_{\nu}-\partial_{\nu}\mathcal A_{\mu}
   +[\mathcal A_{\mu},\mathcal A_{\nu}]
  =-\,P_{\mathfrak h}\bigl([e_{\mu},e_{\nu}]\bigr)
\end{equation}
is, by the Maurer--Cartan structure equation, the projection onto
$\mathfrak h$ of the bracket of two coset directions. This is the gauge
field whose kinetic term we now generate radiatively.

\begin{remark}[Corrected massive multiplet]\label{rem:app-ym:multiplet}
The fluctuations that couple minimally to $\mathcal A_{\mu}$ are the
\emph{coset} directions $\mathfrak m=\mathbf{16}_{-3}\oplus
\overline{\mathbf{16}}_{+3}$, a single real-irreducible $H$-module of
complex type (commutant $\mathbb C$; see
\autoref{sec:foundation:tangent}). It contains \emph{no} $\SOTen$
singlet. Earlier drafts decomposed the massive sector as
$\mathbf{10}_{+2}\oplus\mathbf{10}_{-2}\oplus\mathbf 4_{0}$, with total
Dynkin index $\tfrac52$; that decomposition is incorrect (it places a
nonexistent singlet $\mathbf 4_{0}$ in $\mathfrak m$ and uses
fundamental rather than spinor reps). The corrected weight is computed
below from $\mathbf{16}\oplus\overline{\mathbf{16}}$.
\end{remark}

Expanding the kinetic action to quadratic order about the vacuum, the
$32$ real coset fluctuations $\xi$ acquire the common radiative mass
$m^{2}=f^{2}\mu^{2}$ (one common eigenvalue by Schur on the irreducible
$\mathfrak m$; \autoref{sec:gravity:spectrum}). Of these, $28$ are
physical massive scalars and $4$ are the gauge/soldering frame modes
eaten by diffeomorphism $+$ local Lorentz. The Euclidean quadratic
action in the composite background reads
\begin{equation}\label{eq:app-ym:Squad}
  S_{m}[\xi]
  =\frac12\int_{X}\xi^{\dagger}\bigl[-D^{2}+m^{2}\bigr]\xi\,d^{4}x,
  \qquad
  D_{\mu}=\partial_{\mu}+\mathcal A_{\mu},
\end{equation}
and integrating out $\xi$ produces the determinant
$[\det(-D^{2}+m^{2})]^{-1/2}$, i.e. the one-loop effective action
\begin{equation}\label{eq:app-ym:W}
  \Gamma_{\mathrm{1L}}
  =\frac12\Tr\log(-D^{2}+m^{2})
  =-\frac12\int_{\Lambda^{-2}}^{\infty}\frac{ds}{s}\,
    e^{-sm^{2}}\,\Tr\,e^{-s(-D^{2})}.
\end{equation}

%----------------------------------------------------------------------%
\subsection{Heat-kernel extraction of the gauge weight}
\label{sec:app-ym:heatkernel}
%----------------------------------------------------------------------%

For small proper time the heat trace expands as
\begin{equation}\label{eq:app-ym:HK}
  \Tr\,e^{-s(-D^{2})}
  =\frac{1}{(4\pi s)^{2}}\sum_{n\ge0}s^{n}\,a_{2n},
\end{equation}
and the $\mathcal F^{2}$ content sits in the Seeley--DeWitt coefficient
$a_{4}$. For a real scalar multiplet in an $H$-representation $R$ with
Dynkin index $T(R)$ normalized by
$\Tr_{R}(T^{a}T^{b})=T(R)\,\delta^{ab}$,
\begin{equation}\label{eq:app-ym:a4}
  a_{4}(R)\supset
  \frac{1}{12}\int_{X}\Tr_{R}\bigl(\mathcal F_{\mu\nu}
  \mathcal F^{\mu\nu}\bigr)\,d^{4}x ,
\end{equation}
the omitted terms being curvature-squared invariants on $X$ suppressed
by powers of $\mu/\Lambda$.

\begin{lemma}[Total gauge weight of the coset multiplet]
\label{lem:app-ym:weight}
The total Dynkin index carried by the massive coset multiplet
$\mathfrak m=\mathbf{16}_{-3}\oplus\overline{\mathbf{16}}_{+3}$ is
\begin{equation}\label{eq:app-ym:Ttot}
  T_{\mathfrak m}
  =T(\mathbf{16})+T(\overline{\mathbf{16}})
  =2+2=4 ,
  \qquad
  \dim\mathfrak m=16+16=32 .
\end{equation}
\end{lemma}

\begin{proof}
$\mathfrak m$ is a single real-irreducible $H$-module of complex type
realized as the chiral spinor pair $\mathbf{16}\oplus
\overline{\mathbf{16}}$ of $\SOTen$
(\autoref{sec:foundation:tangent}). The Dynkin index is additive over a
direct sum, and conjugation preserves it, so
$T(\overline{\mathbf{16}})=T(\mathbf{16})=2$ from
\eqref{eq:app-ym:casimirs}. Hence $T_{\mathfrak m}=4$ and
$\dim\mathfrak m=32$. (The $\UOne$ charges $\pm3$ do not enter the
$\SOTen$ trace normalization; they contribute only to the Abelian
factor, treated separately.) The $4$ soldering frame modes are eaten
gauge degrees of freedom carrying no net dynamical $\SOTen$ index, so the
gauge-coupling sum is the same whether weighted by the full $32$ or by
the $28$ physical scalars: the $\SOTen$ trace of any $H$-singlet eaten
combination vanishes. We therefore use the $H$-irreducible weight
$T_{\mathfrak m}=4$. This corrects the legacy value $\tfrac52$ noted in
\autoref{rem:app-ym:multiplet}.
\end{proof}

Inserting \eqref{eq:app-ym:a4} into \eqref{eq:app-ym:W} and using
\autoref{lem:app-ym:weight},
\begin{equation}\label{eq:app-ym:insert}
  \Gamma_{\mathrm{1L}}
  \supset
  -\frac12\int_{\Lambda^{-2}}^{\infty}\frac{ds}{s}\,e^{-sm^{2}}
   \frac{1}{(4\pi)^{2}}\,\frac{s^{2}}{12}\,T_{\mathfrak m}
   \int_{X}\Tr\bigl(\mathcal F_{\mu\nu}\mathcal F^{\mu\nu}\bigr)\,d^{4}x .
\end{equation}
The proper-time integral is logarithmically divergent: with
$\int_{\Lambda^{-2}}^{\infty}ds\,s\,e^{-sm^{2}}/s
=\int_{\Lambda^{-2}}^{\infty}ds\,e^{-sm^{2}}
\;\to\;m^{-2}$ at the upper end and the cutoff supplying
$\log(\Lambda^{2}/m^{2})$ in the $s^{2}\!\cdot s^{-1}$ measure, one
obtains
\begin{equation}\label{eq:app-ym:log}
  \Gamma_{\mathrm{1L}}
  \supset
  -\frac14\,\frac{T_{\mathfrak m}}{24\pi^{2}}\,
   \log\frac{\Lambda^{2}}{m^{2}}
   \int_{X}\Tr\bigl(\mathcal F_{\mu\nu}\mathcal F^{\mu\nu}\bigr)\,d^{4}x .
\end{equation}

%----------------------------------------------------------------------%
\subsection{The induced coupling}
\label{sec:app-ym:coupling}
%----------------------------------------------------------------------%

\begin{theorem}[Induced Yang--Mills coefficient]
\label{thm:app-ym:coupling}
Integrating out the massive clock fields induces the local Yang--Mills
term
\begin{equation}\label{eq:app-ym:YM}
  S_{\mathrm{YM}}
  =-\frac{1}{4\ghat^{2}}\int_{X}
    \Tr\bigl(\mathcal F_{\mu\nu}\mathcal F^{\mu\nu}\bigr)\,d^{4}x ,
\end{equation}
with the one-loop coefficient
\begin{equation}\label{eq:app-ym:ghat}
  \ghat^{-2}
  =\frac{T_{\mathfrak m}}{24\pi^{2}}\,f^{2}\,
   \log\frac{\Lambda^{2}}{\mu^{2}}
  =\frac{4}{24\pi^{2}}\,f^{2}\,\log\frac{\Lambda^{2}}{\mu^{2}}
  =\frac{f^{2}}{6\pi^{2}}\,\log\frac{\Lambda^{2}}{\mu^{2}}
  >0 ,
\end{equation}
where $m^{2}=f^{2}\mu^{2}$ has been restored and $f^{2}$ extracted by
the canonical normalization of the kinetic term.
\end{theorem}

\begin{proof}
Match \eqref{eq:app-ym:log} to \eqref{eq:app-ym:YM}: the coefficient of
$-\tfrac14\Tr(\mathcal F^{2})$ in $\Gamma_{\mathrm{1L}}$ is
$\ghat^{-2}=T_{\mathfrak m}\log(\Lambda^{2}/m^{2})/(24\pi^{2})$ in units
where the $\sigma$-model is canonically normalized; restoring the decay
constant $f$ that sets the overall scale of the Maurer--Cartan kinetic
term multiplies this by $f^{2}$, and $\log(\Lambda^{2}/m^{2})
=\log(\Lambda^{2}/\mu^{2})+\log(1/f^{2})\to\log(\Lambda^{2}/\mu^{2})$ to
leading log. With $T_{\mathfrak m}=4$ from \autoref{lem:app-ym:weight},
$\ghat^{-2}=4f^{2}\log(\Lambda^{2}/\mu^{2})/(24\pi^{2})
=f^{2}\log(\Lambda^{2}/\mu^{2})/(6\pi^{2})$. Positivity follows from
$\Lambda>\mu$, $f^{2}>0$.
\end{proof}

\begin{remark}[Comparison with the legacy coefficient]
\label{rem:app-ym:legacy}
The corrected prefactor $T_{\mathfrak m}/(24\pi^{2})=4/(24\pi^{2})
=1/(6\pi^{2})$ replaces the legacy
$\tfrac52\cdot 1/(24\pi^{2})=5/(48\pi^{2})$, which used the spurious
$\mathbf{10}_{\pm2}\oplus\mathbf 4_{0}$ decomposition with index
$\tfrac52$. The corrected value is larger by a factor
$T_{\mathfrak m}/(5/2)=4/(5/2)=8/5$. The physical content -- a
positive, logarithmically running induced gauge coupling -- is unchanged;
only the numerical prefactor is fixed to the spinor-pair value.
\end{remark}

\begin{remark}[Normalization convention]\label{rem:app-ym:norm}
$\Tr(\mathcal F_{\mu\nu}\mathcal F^{\mu\nu})$ is normalized in the
fundamental $\mathbf{10}$ of $\SOTen$, $\Tr_{\mathbf{10}}(T^{a}T^{b})
=T(\mathbf{10})\,\delta^{ab}$ with $T(\mathbf{10})=1$. Any rescaling of
this trace rescales $\ghat^{2}$ uniformly and drops out of all
physical (renormalization-group) statements below.
\end{remark}

%----------------------------------------------------------------------%
\subsection{Feeding the asymptotic-freedom slope $b_{0}$}
\label{sec:app-ym:b0}
%----------------------------------------------------------------------%

The coefficient \eqref{eq:app-ym:ghat} sets the boundary value of the
gauge coupling at the breaking scale; its scale dependence is governed
by the full one-loop $\beta$-function of the FRG sector
(\autoref{sec:frg:gauge}). With the canonical dimensionless coupling
$\ghat^{2}$ the locked form is
\begin{equation}\label{eq:app-ym:beta}
  \beta_{\ghat^{2}}
  =2\,\ghat^{2}+\frac{b_{0}}{48\pi^{2}}\,\ghat^{4},
  \qquad
  b_{0}=\frac{11}{3}\,C_{A}-\frac{4}{3}\,T\,n_{F}
        -\frac16\,T\,n_{S} ,
\end{equation}
where the canonical $+2\,\ghat^{2}$ term enforces a \emph{Gaussian} UV
fixed point $\ghat^{2}_{\star}=0$: the gauge sector is asymptotically
\emph{free}, and there is no interacting gauge fixed point.

\begin{theorem}[Slope from the corrected matter content]
\label{thm:app-ym:b0}
With the locked Casimirs \eqref{eq:app-ym:casimirs}, the slope is
$b_{0}=26$ for one $\SOTen$ generation ($n_{F}=1$, $n_{S}=2$) and
$b_{0}=12$ for the three-generation theory of \autoref{sec:sm:generations}
($n_{F}=3$, $n_{S}=28$).
\end{theorem}

\begin{proof}
\emph{One generation.} Insert $C_{A}=8$, $T=T(\mathbf{16})=2$, $n_{F}=1$,
and $n_{S}=2$ (minimal GUT scalars at the matching scale):
\begin{equation}\label{eq:app-ym:b0one}
  b_{0}(1)
  =\frac{11}{3}\cdot 8-\frac43(2)(1)-\frac16(2)(2)
  =\frac{78}{3}=26 .
\end{equation}
\emph{Three generations.} With three chiral $\mathbf{16}$ representations
($n_{F}=3$, not $3\times16$) and $n_{S}=28$ massive coset scalars,
\begin{equation}\label{eq:app-ym:b0lin}
  b_{0}(3)
  =\frac{11}{3}(8)-\frac43(2)(3)-\frac16(2)(28)
  =\frac{88-24-28}{3}=12 .
\end{equation}
Hence $b_{0}(3\,\mathrm{gen})=12$, the value used everywhere downstream.
\end{proof}

\begin{remark}[Why the slope does not change the UV behavior]
\label{rem:app-ym:freedom}
The magnitude of $b_{0}$ controls the \emph{rate} of running but not the
existence of the UV fixed point: the canonical $+2\,\ghat^{2}$ term in
\eqref{eq:app-ym:beta} produces $\ghat^{2}_{\star}=0$. In either case the
gauge sector is asymptotically free, and the induced coefficient
\eqref{eq:app-ym:ghat} supplies a finite, positive boundary value at the
breaking scale.
\end{remark}

This closes the chain from the Maurer--Cartan curvature
\eqref{eq:app-ym:curv}, through the corrected coset weight
$T_{\mathfrak m}=4$ (\autoref{lem:app-ym:weight}), to the induced
coupling \eqref{eq:app-ym:ghat} and the asymptotic-freedom slope
$b_{0}=12$ for three generations (\autoref{thm:app-ym:b0}) carried into
the FRG analysis of \autoref{sec:frg:gauge}.
